\chapter{The opponent: three measured strategies}
\label{chap-adversaire}

\section{A pure function}

The artificial opponent comes down to one function (\cref{lst-ia}): given a level, the public
knowledge of the board fired at and a random generator, it designates the next cell to aim
at.

\begin{lstlisting}[caption={The whole opponent (\fichier{include/battleship/ai.h}).},
                   label={lst-ia}]
bs_status bs_ai_choose_target(bs_difficulty difficulty,
                              const bs_knowledge *knowledge,
                              bs_rng *rng, bs_coord *target);
\end{lstlisting}

This function \textbf{keeps no memory} from one call to the next. \enquote{A ship is
    wounded} is not a piece of information to remember: it can be read on the board, where a
\code{HIT} cell designates a deck that is hit on a ship still afloat.

\begin{attention}
    In version~1, the pursuit of a hit ship relied on a global variable. After a hit, the
    computer drew one of the four neighbours at random; if it was off the board or already
    played, it \emph{gave up the pursuit} and went back to random fire, forgetting the wounded
    ship. Worse, the cell drawn was read from the array \emph{before} checking that it was on
    the board: an out-of-bounds read, undefined behaviour in C.
\end{attention}

A function without memory cannot fall out of step with the game. It also lends itself to
every experiment: any position can be submitted to it, it can be made to play against
itself, or used to advise the player.

\section{Three levels, one more skill at each step}

\begin{table}[ht]
    \centering
    \begin{tabularx}{\textwidth}{@{}l>{\raggedright\arraybackslash}Xccc@{}}
        \toprule
        \textbf{Level}                       & \textbf{Strategy}                                                     &
        \rotatebox{60}{\small never repeats} & \rotatebox{60}{\small deduces}                                        &
        \rotatebox{60}{\small finishes, anticipates}                                                                             \\
        \midrule
        Easy                                 & Fires at random among the cells not yet fired at.                     &
        $\bullet$                            &                                                                       &           \\
        Medium                               & Fires at random, but never at a cell that the rules prove empty.      &
        $\bullet$                            & $\bullet$                                                             &           \\
        Hard                                 & Finishes off wounded ships; when hunting, aims at the likeliest cell. &
        $\bullet$                            & $\bullet$                                                             & $\bullet$ \\
        \bottomrule
    \end{tabularx}
    \caption{The three levels of the opponent.}
    \label{tab-niveaux}
\end{table}

The easy level ignores the deductions of \cref{sec-deductions}: it can fire at a cell marked
with a flag. The medium level uses them, and that is its only difference. The hard level
adds reasoning about the ships that remain.

\section{The hard level: a density map}

\subsection{Principle}

The hard level asks itself a simple question: \emph{among all the ways of placing the ships
    still afloat that are compatible with what I know, how many go through this cell?} The cell
that the most placements go through is the likeliest.

Let $K$ be the current knowledge, $n_\ell$ the number of ships of length $\ell$ still afloat
and $P_\ell(K)$ the set of consistent positions for such a ship. The score of a cell $c$
that is still unknown is
\begin{equation}
    s(c) \;=\; \sum_{\ell = 1}^{4} n_\ell \cdot
    \bigl|\{\, p \in P_\ell(K) \;:\; c \in p \,\}\bigr| .
    \label{eq-densite}
\end{equation}

A position $p$ is \emph{consistent} with $K$ when the following three conditions hold:

\begin{enumerate}
    \item each of its cells is either unknown or a hit deck of a ship afloat;
    \item it is not made \emph{only} of hit decks --- that ship would already be sunk;
    \item no hit deck adjoins it without belonging to it --- two ships do not touch.
\end{enumerate}

\Cref{fig-densite} shows the map obtained on a blank board. A corner is covered by only 16
placements, the centre by 36: without any rule telling it so, the computer opens the game in
the centre, where the large ships have the most ways of being.

\begin{figure}[ht]
    \centering
    \begin{tikzpicture}[scale=0.58]
        % Real scores computed by bs_ai_density() on a blank board (symmetric).
        \foreach \v [count=\i from 0] in {
                16,22,25,26,26,26,26,25,22,16,
                22,28,31,32,32,32,32,31,28,22,
                25,31,34,35,35,35,35,34,31,25,
                26,32,35,36,36,36,36,35,32,26,
                26,32,35,36,36,36,36,35,32,26,
                26,32,35,36,36,36,36,35,32,26,
                26,32,35,36,36,36,36,35,32,26,
                25,31,34,35,35,35,35,34,31,25,
                22,28,31,32,32,32,32,31,28,22,
                16,22,25,26,26,26,26,25,22,16} {
                \pgfmathtruncatemacro{\col}{mod(\i,10)}
                \pgfmathtruncatemacro{\row}{9 - int(\i/10)}
                \pgfmathtruncatemacro{\teinte}{(\v - 16) * 5}
                \fill[ocean!\teinte!ecume] (\col,\row) rectangle ++(1,1);
                \ifnum\v>29
                    \node[font=\sffamily\scriptsize, text=white] at (\col+0.5,\row+0.5) {\v};
                \else
                    \node[font=\sffamily\scriptsize, text=marine] at (\col+0.5,\row+0.5) {\v};
                \fi
            }
        \plateau{10}
        \foreach \l [count=\i from 0] in {A,B,C,D,E,F,G,H,I,J}
        \node[font=\sffamily\scriptsize, text=ardoise] at (\i+0.5,10.45) {\l};
        \foreach \n in {1,...,10}
        \node[font=\sffamily\scriptsize, text=ardoise] at (-0.45,10.5-\n) {\n};
    \end{tikzpicture}
    \caption{Density map of a blank board: number of placements of the fleet covering each
        cell (values produced by \code{bs\_ai\_density()}).}
    \label{fig-densite}
\end{figure}

\subsection{Hunting and finishing off with the same code}

As long as no ship is wounded, all the consistent placements count: this is the \emph{hunt}.
As soon as a \code{HIT} cell exists, only the placements that \emph{explain} at least one hit
deck are kept, weighted by the number of decks explained. The same formula then becomes a
procedure for \emph{finishing off}:

\begin{itemize}
    \item around an isolated hit deck, the four neighbours receive a score, all the higher
          as long ships can extend in their direction;
    \item as soon as two aligned decks are hit, condition~3 eliminates every perpendicular
          position: only the two cells extending the line keep a score.
\end{itemize}

No special case is programmed for this; the behaviour follows from the three consistency
conditions (\cref{alg-densite}).

\begin{algorithm}[ht]
    \caption{Choice of a target at the hard level}
    \label{alg-densite}
    \KwIn{the knowledge $K$, a random generator $g$}
    \KwOut{a cell to aim at}
    $s(c) \gets 0$ for every cell $c$\;
    $\mathit{finish} \gets$ there is a \textsc{hit} cell in $K$\;
    \ForEach{length $\ell$ such that $n_\ell > 0$}{
        \ForEach{position $p$ of a ship of length $\ell$ on the board}{
            \lIf{$p$ is not consistent with $K$}{skip to the next one}
            $h \gets$ number of \textsc{hit} cells of $p$\;
            \lIf{$\mathit{finish}$ and $h = 0$}{skip to the next one}
            $w \gets n_\ell \times h$ if $\mathit{finish}$, $n_\ell$ otherwise\;
            \lForEach{unknown cell $c$ of $p$}{$s(c) \gets s(c) + w$}
        }
    }
    \Return a cell of maximal score, drawn at random with $g$ in case of a tie\;
\end{algorithm}

The cost is negligible: four lengths, two orientations, a hundred origins, ships of four
cells at most --- a few thousand operations per shot. Drawing at random among cells of equal
score keeps the computer from always playing the same opening.

\section{Measuring rather than assuming}

\subsection{Protocol}

The program \code{battleship-cli simulate} makes the strategies play without an interface.
Two experiments are run, entirely determined by a seed:

\begin{description}[style=nextline, leftmargin=1.2em]
    \item[Solo] One strategy fires at a random fleet until it sinks it. The shots are
          counted: at least 20 are needed, at most 100. This measurement isolates the quality of
          the strategy.
    \item[Duel] Two strategies face each other with the real rule \enquote{a hit plays
              again}. Each pairing is played with the starting side alternating, to cancel out the
          advantage of the first shot.
\end{description}

All the figures of this chapter come from the command
\code{battleship-cli simulate -{}-games 10000 -{}-seed 1} and can be reproduced exactly.

\subsection{Results}

The levels are clearly ordered (\cref{tab-solo,fig-histogramme}): it takes on average
\num{96.2} shots at the easy level, \num{65.8} at the medium level and \num{55.3} at the hard
level. In a duel (\cref{tab-duels}), the hard level beats the medium one in
\SI{87.5}{\percent} of the games, and the medium level wins \num{9999} of its \num{10000}
games against the easy one. The diagonal, close to \SI{50}{\percent}, confirms that
alternating the sides does neutralise the advantage of the first shot.

\begin{table}[H]
    \centering
    \begin{tabular}{@{}l S[table-format=2.2] S[table-format=1.2] S[table-format=2.0]
            S[table-format=2.0] S[table-format=3.0] S[table-format=3.0]@{}}
        \toprule
        \textbf{Level}    & {\textbf{Mean}}                         & {\textbf{Std dev.}} & {\textbf{Min.}} &
        {\textbf{Median}} & {\textbf{9\textsuperscript{th} decile}} & {\textbf{Max.}}                                        \\
        \midrule
        Easy              & 96.19                                   & 4.10                & 69              & 97 & 100 & 100 \\
        Medium            & 65.83                                   & 6.07                & 44              & 66 & 74  & 85  \\
        Hard              & 55.27                                   & 6.80                & 28              & 56 & 64  & 79  \\
        \bottomrule
    \end{tabular}
    \caption{Number of shots needed to sink a random fleet (\num{10000} games per level).}
    \label{tab-solo}
\end{table}

\begin{figure}[H]
    \centering
    \begin{tikzpicture}
        \begin{axis}[
                width=0.96\textwidth, height=6.4cm,
                xlabel={Number of shots to sink the fleet},
                ylabel={Games (out of \num{10000})},
                xmin=25, xmax=101, ymin=0,
                xtick={30,40,...,100}, ytick distance=500,
                scaled y ticks=false,
                yticklabel={\pgfmathprintnumber[fixed, 1000 sep={,}]{\tick}},
                axis lines=left, axis line style={draw=ardoise},
                tick style={draw=ardoise}, tick label style={font=\sffamily\footnotesize},
                label style={font=\sffamily\small},
                ymajorgrids, grid style={draw=ardoise!20},
                legend style={draw=none, fill=none, font=\sffamily\small, at={(0.03,0.97)},
                        anchor=north west, legend cell align=left},
                const plot mark mid, area style, enlarge x limits=false,
            ]
            \addplot[draw=vert, fill=vert, fill opacity=0.30, line width=0.7pt]
            table[x=shots, y=hard] {data/benchmark.dat} \closedcycle;
            \addlegendentry{Hard}
            \addplot[draw=ocean, fill=ocean, fill opacity=0.30, line width=0.7pt]
            table[x=shots, y=medium] {data/benchmark.dat} \closedcycle;
            \addlegendentry{Medium}
            \addplot[draw=corail, fill=corail, fill opacity=0.30, line width=0.7pt]
            table[x=shots, y=easy] {data/benchmark.dat} \closedcycle;
            \addlegendentry{Easy}
        \end{axis}
    \end{tikzpicture}
    \caption{Distribution of the number of shots by level. The easy level runs into the
        ceiling of 100 shots; the other two form well-separated bell curves.}
    \label{fig-histogramme}
\end{figure}

\begin{table}[H]
    \centering
    \begin{tabular}{@{}l S[table-format=3.1] S[table-format=3.1] S[table-format=2.1]@{}}
        \toprule
                        & {\textbf{against Easy}} & {\textbf{against Medium}} & {\textbf{against Hard}} \\
        \midrule
        \textbf{Easy}   & 49.6                    & 0.0                       & 0.0                     \\
        \textbf{Medium} & 100.0                   & 49.7                      & 12.0                    \\
        \textbf{Hard}   & 100.0                   & 87.5                      & 50.0                    \\
        \bottomrule
    \end{tabular}
    \caption{Duels: percentage of games won by the level of the row (\num{10000} games per
        cell).}
    \label{tab-duels}
\end{table}

\subsection{How the levels were chosen}

The three levels do not come from an intuition but from a comparison. Four candidate
strategies were measured during development (\cref{tab-candidats}).

\begin{table}[H]
    \centering
    \begin{tabularx}{\textwidth}{@{}>{\raggedright\arraybackslash}X S[table-format=2.1]
            S[table-format=+2.1] l@{}}
        \toprule
        \textbf{Strategy}                              & {\textbf{Shots}} & {\textbf{Gain}} & \textbf{Kept as} \\
        \midrule
        A. Pure chance                                 & 96.2             &                 & Easy             \\
        B. A + respect for the deductions              & 65.8             & -30.4           & Medium           \\
        C. B + classic pursuit of hit ships            & 58.2             & -7.6            & ---              \\
        D. B + density map (hunting and finishing off) & 55.3             & -2.9            & Hard             \\
        \bottomrule
    \end{tabularx}
    \caption{Candidate strategies. The gain is expressed relative to the previous row;
        strategy~C was measured over \num{2000} games.}
    \label{tab-candidats}
\end{table}

The lesson of this table is counter-intuitive. With the no-contact rule, the skill that
matters most is neither the pursuit of hit ships nor the computation of probabilities: it is
the simple fact of \textbf{not firing where the rules prove there is nothing}. On its own,
this discipline saves thirty shots. The classic pursuit adds fewer than eight, and the
density map three more.

The levels were therefore spaced according to these measurements: strategy~C, too close
to~D, would not have made a noticeable intermediate level.

\begin{retenir}[A floor that is hard to lower]
    Why does the best level still need 55 shots for 20 cells to hit? Because of the four ships
    of a single cell. No reasoning can \emph{locate} them: once the large ships are sunk,
    nearly all the remaining cells have to be explored. The density map gets close to this
    floor; it cannot remove it.
\end{retenir}
